% Generated draft: requires mathematical and source review.
\documentclass[11pt]{article}
\usepackage[margin=25mm]{geometry}
\usepackage{fontspec}
\setmainfont{Latin Modern Roman}
\usepackage{amsmath,amssymb,mathtools,unicode-math}
\setmathfont{Latin Modern Math}
\usepackage{enumitem,needspace,xurl,hyperref}
\setlength{\parindent}{0pt}
\setlength{\parskip}{5pt}
\setlength{\emergencystretch}{3em}
\newcommand{\N}{\mathbb{N}}
\newcommand{\Z}{\mathbb{Z}}
\newcommand{\Q}{\mathbb{Q}}
\newcommand{\R}{\mathbb{R}}
\newcommand{\C}{\mathbb{C}}
\newcommand{\F}{\mathbb{F}}
\newcommand{\A}{\mathbb{A}}
\newcommand{\PP}{\mathbb{P}}
\newcommand{\E}{\mathbb{E}}
\newcommand{\eps}{\varepsilon}
\newcommand{\dd}{\,\mathrm{d}}
\begin{document}
\section*{UnsolvedMath ID: 800242}
\textbf{Source problem code:} catalogue entry
\subsection{Definitions and mathematical statement}
For a positive integer \(n\), let \(D\subseteq\C^n\) be a domain, meaning a nonempty connected open set, and let \(dV\) denote Euclidean volume measure. Its Bergman space is
\[
A^2(D)=\left\{f:D\to\C:\ f\text{ is holomorphic and }\int_D|f|^2\,dV<\infty\right\}.
\]
The Bergman kernel \(K_D(z,w)\) is holomorphic in \(z\), antiholomorphic in \(w\), satisfies \(K_D(\cdot,w)\in A^2(D)\), and has the reproducing property
\[
f(z)=\int_D K_D(z,w)f(w)\,dV(w)
\qquad(f\in A^2(D),\ z\in D).
\]
The source asks for a characterization of the domains \(D\) for which
\[
\forall z,w\in D,\qquad K_D(z,w)\ne0.
\]
This condition concerns all pairs of interior points, including distinct points; it imposes no boundary condition. If \(A^2(D)=\{0\}\), then \(K_D\equiv0\), so such a domain does not satisfy the condition.

The source does not specify the dimension or restrict the domains to a bounded, smooth, pseudoconvex, or other particular class. The definitions above provide a general framework, but the intended scope of the requested characterization remains unspecified. The source relates the question to the Lu Qi-Keng conjecture without stating that conjecture precisely.
\subsection{Short English statement}
Characterize the domains whose Bergman kernel has no zeros anywhere on \(D\times D\). The supplied source leaves the dimension and intended class of domains unspecified and describes the question as related to the Lu Qi-Keng conjecture.
\subsection{Sources}
\begin{itemize}
\item Local UnsolvedMath registry, permanent ID \(800242\), source code \texttt{catalogue entry}, title ``Several Complex Variables Problem List --- Problem 39.'' The supplied registry attributes the problem to Buzzard, Fornaess, and Krantz, with proposed year 1995, and labels it ``open.'' This status is supplied metadata, not an independently verified mathematical or literature assessment.
\item Buzzard, Fornaess, and Krantz, \emph{The Several Complex Variables Problem List} (1995), Problem 39. Supplied source URL: \url{https://arxiv.org/abs/math/9509202}. The record reports access on 2026-10-16 and extraction by \texttt{arxiv\_latex}; these are supplied provenance details, not a claim of independent access or verification.
\end{itemize}
\end{document}
